\documentclass[aspectratio=169]{beamer}
\usepackage{tikz}
\usepackage{graphicx}
\usepackage{listings}
\usepackage{xcolor}
\usepackage{hyperref}
\usepackage[spanish,es-noquoting]{babel}
\usepackage[round]{natbib}

\lstset{
  basicstyle=\ttfamily\small,
  keywordstyle=\color{blue},
  stringstyle=\color{red},
  commentstyle=\color{gray},
  showstringspaces=false
}
% Minimal look
\usetheme{Warsaw}
%\setbeamertemplate{headline}
\setbeamertemplate{footline}{
  \leavevmode%
  \hbox{%
    % Left: Section name
    \begin{beamercolorbox}[wd=.33\paperwidth,ht=2.25ex,dp=1ex,left]{author in head/foot}%
      \hspace*{2ex}\insertsection
    \end{beamercolorbox}%
    % Center: Presentation title
    \begin{beamercolorbox}[wd=.34\paperwidth,ht=2.25ex,dp=1ex,center]{title in head/foot}%
      \insertshorttitle
    \end{beamercolorbox}%
    % Right: Slide numbers
    \begin{beamercolorbox}[wd=.33\paperwidth,ht=2.25ex,dp=1ex,right]{date in head/foot}%
      \insertframenumber{} / \inserttotalframenumber\hspace*{2ex}
    \end{beamercolorbox}%
  }%
  \vskip0pt%
}

\usecolortheme{default}
\usefonttheme{serif}
\graphicspath{{../}{./}}
\setbeamertemplate{caption}[numbered]

% Title info
\title[Módulo 4: Explorando la Creati]{Módulo 4: Explorando la Creatividad Visual}
\subtitle{Generación de Imágenes y Multimedia con IA}
\author{Ricardo Reyes-Jimenez}
\institute{IA para Todos}
\date{\today}



\begin{document}

% Title slide
\begin{frame}
  \titlepage
\end{frame}

% Reset background for normal slides
\usebackgroundtemplate{}



\begin{frame}[fragile]{Objetivo de la Sesión}


\textbf{Al finalizar este módulo, los participantes serán capaces de:}

1. Comprender cómo la IA puede crear imágenes y diseños desde cero.
2. Usar la IA para tareas visuales cotidianas (presentaciones, redes sociales, regalos).
3. Conocer herramientas accesibles para dar vida a sus ideas.

\end{frame}



\begin{frame}[fragile]{1. El Fin de las Imágenes Genéricas}


\textbf{Ya no dependemos de Google Imágenes. Si lo puedes imaginar, lo puedes crear.}

\begin{itemize}
  \item \textbf{Caso Práctico:} Un padre organizando una fiesta de cumpleaños temática para su hijo (ej. ``Dinosaurios en el espacio''), y necesita invitaciones únicas.
  \item \textbf{Uso de IA (Midjourney/DALL-E):} ``Crea una ilustración colorida estilo Pixar de un T-Rex astronauta sosteniendo un globo, flotando en el espacio con planetas de fondo. Deja espacio libre a la derecha para añadir texto.''
\end{itemize}
\textbf{Resultado:} Diseños únicos en minutos sin necesidad de ser un experto en Photoshop.

\end{frame}



\begin{frame}[fragile]{2. Rediseñando Espacios (Arquitectos y Decoradores)}


\textbf{La IA te ayuda a visualizar cambios antes de gastar un solo centavo.}

\begin{itemize}
  \item \textbf{Caso Práctico:} Quieres remodelar tu sala, pero no sabes si un estilo industrial o minimalista quedará bien.
  \item \textbf{Uso de IA:} Usando herramientas de IA para diseño de interiores, subes una foto de tu sala actual y pides: ``Transforma esta habitación en un diseño moderno de mediados de siglo (mid-century modern) con plantas de interior y luz natural cálida.''
\end{itemize}
\textbf{Resultado:} Múltiples ideas de decoración instantáneas para inspirarte o mostrárselas a tu contratista.

\end{frame}



\begin{frame}[fragile]{3. Del Texto a la Presentación}


\textbf{No más horas perdidas ajustando cajitas de texto en PowerPoint.}

\begin{itemize}
  \item \textbf{Caso Práctico:} Un estudiante o profesor necesita armar una presentación rápida sobre el cambio climático, pero tiene poco tiempo.
  \item \textbf{Uso de IA (Gamma/Tome):} Provees un esquema o un texto básico a la IA y le pides: ``Crea una presentación de 5 diapositivas sobre el impacto del cambio climático en los océanos. Usa un diseño limpio y moderno con tonos azules.''
\end{itemize}
\textbf{Resultado:} Un borrador visualmente atractivo listo para ser afinado.

\end{frame}



\begin{frame}[fragile]{Recapitulando el Módulo}


\textbf{Hoy hemos aprendido:}
\begin{itemize}
  \item Que la generación visual es una herramienta para potenciar nuestra imaginación, no solo para artistas.
  \item Cómo describir texturas, estilos y detalles visuales a la IA.
  \item A visualizar espacios, eventos y presentaciones de forma rápida y económica.
\end{itemize}
\textbf{¡Sigan aprendiendo, sigan trabajando y sigan disfrutando de la Inteligencia Artificial!}
\end{frame}




\begin{frame}[allowframebreaks]{References}
  \bibliographystyle{apalike}
  \bibliography{references}
\end{frame}


\end{document}